\documentclass[aps,prl,twocolumn,superscriptaddress,floatfix]{revtex4-2}
\usepackage{amsmath,amssymb,amsfonts}
\usepackage{bm}
\usepackage{hyperref}
\usepackage{xcolor}
\usepackage{booktabs}

\definecolor{bctnavy}{RGB}{10,36,68}
\definecolor{bctblue}{RGB}{13,71,161}
\definecolor{bctgreen}{RGB}{27,94,32}

\begin{document}

\title{BCT Letter 125: The Faraday Test\\
\large A Decisive Experimental Discrimination Between OHC Vacuum Coupling\\
\large and Electromagnetic Substrate Theories of Morphic Resonance}

\author{Michel Robert Cabri\'{e}}
\affiliation{Independent Researcher, Victoria, Australia}
\email{ZeroFreeParameters@gmail.com}
\date{\today}

\begin{abstract}
BCT Letter~79 identified the Octet-Hopfion Condensate (OHC) as the
physical substrate of morphic fields~\cite{Sheldrake1981}, deriving
Sheldrake's central empirical claim from first principles via the BCT
Repetition Theorem. One of five resulting predictions is unambiguous and
immediately testable with minimal apparatus: OHC coupling is
charge-independent with coupling constant $\alpha_0^2 \approx 5.5 \times
10^{-5}$, entirely independent of electromagnetic shielding. A Faraday
cage that eliminates all electromagnetic radiation has
\emph{no effect} on OHC phase propagation.
Prediction~\#166: morphic resonance effect sizes are equal inside and
outside a properly constructed Faraday cage, to within experimental
precision. We present a complete experimental protocol
executable by independent researchers with modest apparatus, designed for
collaboration with Prof.\ Sheldrake's existing experimental programme.
The experiment is decisive: it distinguishes the OHC substrate theory from
every electromagnetic-substrate theory of morphic resonance in a single
clean test. Zero free parameters.
\end{abstract}

\maketitle

%------------------------------------------------------------
\section{The Theoretical Prediction}
%------------------------------------------------------------

The BCT Repetition Theorem~\cite{BCT_L79} establishes that every
topological event in the OHC vacuum imprints a permanent phase pattern
propagating at $v_\mathrm{ph} = 6.78c$ with coupling constant
$\alpha_0^2 \approx 5.5 \times 10^{-5}$. The coupling is charge-independent
because it arises from Hopf winding number conservation in
$\pi_3(S^2) = \mathbb{Z}$, not from electromagnetic interaction.

A Faraday cage eliminates propagating electromagnetic fields by
charge redistribution on a conducting surface. This mechanism operates
exclusively on fields that couple to electric charge. The OHC phase
field does not couple to electric charge. Therefore:

\noindent\fbox{\parbox{\dimexpr\columnwidth-2\fboxsep-2\fboxrule\relax}{%
\textbf{Prediction~\#166 (The Faraday Test):}
Morphic resonance effect sizes, measured using any of Sheldrake's
validated protocols, are equal inside and outside a properly
constructed Faraday cage: $R_\mathrm{shielded}/R_\mathrm{unshielded}
= 1.000 \pm \epsilon$, where $\epsilon$ is experimental noise.
\textbf{Zero free parameters.}}}

This prediction distinguishes BCT from every theory in which morphic
effects are mediated by or coupled to electromagnetic fields.

%------------------------------------------------------------
\section{Why This Experiment Is Decisive}
%------------------------------------------------------------

Four classes of competing theory make different predictions:

\begin{enumerate}
\item \textbf{BCT/OHC substrate:} $R_\mathrm{shielded}
= R_\mathrm{unshielded}$. No shielding effect at any frequency.

\item \textbf{Electromagnetic carrier theories}
(any model in which information is carried by EM radiation): complete
suppression of morphic effects inside the Faraday cage.

\item \textbf{Quantum entanglement theories}
(decoherence-based): partial suppression depending on shielding
effectiveness and frequency.

\item \textbf{Null hypothesis}
(no morphic effect exists): equal results in both conditions ---
indistinguishable from BCT in the Faraday test, but distinguishable
from BCT by the pre-existing Sheldrake effect-size literature.
\end{enumerate}

If a morphic effect exists \emph{and} is unaffected by shielding, BCT
is the only framework that predicted this outcome from first principles
before the experiment was run.

%------------------------------------------------------------
\section{Experimental Protocol}
%------------------------------------------------------------

We adopt Sheldrake's validated \emph{staring detection} paradigm
\cite{Sheldrake2000} as the primary test, with the \emph{telephone
telepathy} protocol \cite{Sheldrake2003} as secondary replication.
Both protocols have published positive effect sizes, established
methodology, and volunteer recruitment pathways.

\subsection{Apparatus}

A Faraday cage suitable for this experiment requires:
\begin{enumerate}
\item Copper mesh (aperture $< 1\,\mathrm{mm}$) on all six faces of
    a room-sized or box-sized enclosure, with overlapping seams
    soldered or clamped.
\item Verified shielding effectiveness $> 40\,\mathrm{dB}$ at
    $10\,\mathrm{MHz}$ (suppression factor $> 100\times$).
\item Wooden or plastic frame --- no ferromagnetic components.
\item Battery-powered lighting and recording equipment inside
    (no penetrating cables).
\end{enumerate}

A shielded room of this specification is available commercially
(cost $\sim$A\$2,000--5,000 for a small enclosure) or constructable
from copper mesh and timber. Shielding effectiveness can be verified
with a signal generator and handheld RF meter.

Critically: the cage need not be perfect. It need only attenuate EM
fields by the factor required for a null result under EM-carrier
theories ($> 99\%$ suppression). BCT predicts zero attenuation of
the morphic effect regardless of cage quality.

\subsection{Staring Detection Protocol}

Following Sheldrake~\cite{Sheldrake2000}:

\begin{enumerate}
\item \textbf{Receiver} sits inside the Faraday cage (shielded
    condition) or in an equivalent room without shielding (unshielded
    condition), facing away from the observation direction.
\item \textbf{Sender} observes the receiver via a one-way video feed
    from a separate room. The video cable is shielded and filtered.
\item In each trial, the sender either stares at the receiver
    or looks away, following a randomised schedule generated by
    sealed envelope (pre-registered).
\item The receiver presses a button to indicate whether they feel
    they are being stared at.
\item \textbf{$N = 200$ trials per condition} (shielded/unshielded),
    matched for time of day, sender-receiver pair, and session length.
\end{enumerate}

\subsection{Primary Outcome Measure}

Hit rate $p$ vs.\ chance $(0.50)$ in each condition.
BCT prediction: $p_\mathrm{shielded} = p_\mathrm{unshielded}$.
EM-carrier prediction: $p_\mathrm{shielded} \approx 0.50$
(chance), $p_\mathrm{unshielded} > 0.50$.

\subsection{Statistical Power}

Sheldrake's published staring detection effect size is
$h \approx 0.10$--$0.15$ (Cohen's $h$, corresponding to
$p \approx 0.55$--$0.57$). For $N = 200$ trials per condition
and $\alpha = 0.05$ (two-tailed), power to detect this effect
is $> 80\%$. For $N = 400$ per condition, power exceeds $95\%$.

The critical test is the \emph{comparison} between conditions:
the two-sample test of $p_\mathrm{shielded}$ vs.\
$p_\mathrm{unshielded}$. Under BCT, the two hit rates are equal.
Under EM-carrier theories, they differ significantly.

%------------------------------------------------------------
\section{Pre-Registration and Blinding}
%------------------------------------------------------------

The protocol should be pre-registered at the Open Science Framework
(\url{https://osf.io}) before data collection. Key pre-registration
items:
\begin{itemize}
\item Primary hypothesis and statistical test.
\item Randomisation procedure for staring/not-staring schedule.
\item Faraday cage construction and verified shielding effectiveness.
\item Sample size and stopping rule.
\item Analysis plan (paired vs.\ independent conditions).
\end{itemize}

Pre-registration eliminates the primary objection to morphic resonance
research: post-hoc analysis. A pre-registered null result for shielding
is as informative as a positive result.

%------------------------------------------------------------
\section{Secondary Protocol: Telephone Telepathy}
%------------------------------------------------------------

Sheldrake's telephone telepathy protocol~\cite{Sheldrake2003}:
caller is selected randomly from four known contacts; receiver
guesses which caller before answering. Effect size $h \approx 0.20$
in published studies.

Replication with receiver inside Faraday cage (mobile phone on
battery, isolated from external networks during the guessing phase,
with call delivered via shielded landline or delayed playback).

BCT prediction: identical effect size inside and outside cage.

%------------------------------------------------------------
\section{Implementation: Victoria, Australia}
%------------------------------------------------------------

This experiment is executable immediately from a domestic or
studio setting:

\begin{enumerate}
\item \textbf{Faraday cage:} A purpose-built copper mesh box
    (1m $\times$ 1m $\times$ 1m, cost $\sim$A\$400--800 in materials)
    provides adequate shielding for RF frequencies. Construction
    requires copper mesh, timber framing, and soldering.
    Alternatively, a commercial shielded tent (Faraday bag for
    rooms, widely available) can be used for initial trials.

\item \textbf{Participants:} Sheldrake's published volunteer recruitment
    methodology is straightforward; participants can be recruited via
    university notice boards, online communities, or direct invitation.
    Minimum: 4 sender-receiver pairs, 100 trials each
    condition = 400 trials total.

\item \textbf{Recording:} Battery-powered laptop inside cage for
    trial recording. External camera feed via optical fibre
    (non-conducting, does not break Faraday shielding).

\item \textbf{Timeline:} Construction and verification (2 weeks);
    participant recruitment (2--4 weeks); data collection (4--8 weeks);
    pre-registration and submission (2 weeks).
    \textbf{Total: approximately 3 months from start to submitted
    results.}
\end{enumerate}

%------------------------------------------------------------
\section{Collaboration Proposal}
%------------------------------------------------------------

This Letter constitutes a formal invitation to collaborative
replication. The experiment is designed to:
\begin{enumerate}
\item Be runnable independently by the BCT programme in Victoria,
    Australia, without institutional affiliation.
\item Be directly compatible with Sheldrake's existing protocols,
    enabling meta-analysis combining Victoria results with any
    future replications by Sheldrake's group.
\item Produce a pre-registered, peer-reviewable dataset regardless
    of outcome.
\end{enumerate}

A positive result (morphic effect unaffected by shielding) would
constitute the first experimental confirmation of a specific BCT
prediction in the biological domain.

A negative result (morphic effect eliminated by shielding) would
require revision of the OHC substrate assignment while leaving the
BCT Repetition Theorem intact for other substrates.

Either outcome advances the field. The experiment is decisive.

%------------------------------------------------------------
\section{Conclusion}
%------------------------------------------------------------

The Faraday Test is the cleanest currently executable discrimination
between the BCT OHC substrate theory and electromagnetic-carrier
theories of morphic resonance. It requires modest apparatus, follows
established protocol, and produces a pre-registered result. BCT
Prediction~\#166 is unambiguous: shielding makes no difference.

\textit{The morphic field, if it exists, propagates through the
quantum vacuum itself. Copper mesh is irrelevant to topology.}

\begin{acknowledgments}
This experiment is proposed in response to correspondence with
Prof.\ Rupert Sheldrake (March 2026), who noted that empirical
research has not kept pace with theoretical development in morphic
resonance. BCT Letter~125 is an attempt to help close that gap.
Dedicated to everyone who has looked at the staring effect data and
thought: \textit{the mechanism matters.}
Zeta Vera (CSSC) for computational and editorial support; and the
Gang Gang Cockatoos (\textit{Callocephalon fimbriatum}), who stare
at Michel from the pond edge with an effect size of approximately
1.000 regardless of electromagnetic shielding.

\noindent{\small\textbf{Patreon:}
\href{https://www.patreon.com/cw/TheBCTSuperfluidLatticeModel}{patreon.com/TheBCTSuperfluidLatticeModel}}\\
{\small\textbf{Archive:}
\href{https://zenodo.org/search?q=cabrie}{zenodo.org/search?q=cabri\'e}\quad
\textbf{ORCID:} 0009-0007-9561-9859}
\end{acknowledgments}

\begin{thebibliography}{9}
\bibitem{Sheldrake1981}
R.~Sheldrake, \textit{A New Science of Life},
Blond \& Briggs, London (1981).

\bibitem{Sheldrake2000}
R.~Sheldrake, \textit{J.\ Sci.\ Explor.}~\textbf{14}, 381 (2000)
--- staring detection protocol.

\bibitem{Sheldrake2003}
R.~Sheldrake and M.~Smart,
\textit{J.\ Parapsychol.}~\textbf{67}, 187 (2003)
--- telephone telepathy protocol.

\bibitem{BCT_L79}
M.~R.~Cabri\'{e}, BCT Letter~79: Morphic Resonance and the OHC
Condensate, Zenodo (2026),
\url{https://zenodo.org/search?q=cabrie}

\bibitem{BCT_GoE}
M.~R.~Cabri\'{e}, \textit{BCT Letter Series}, Zenodo (2026),
\url{https://zenodo.org/search?q=cabrie}
\end{thebibliography}

\end{document}
